\documentclass[10pt,a4paper]{amsart}
\usepackage[margin=19mm]{geometry}
\usepackage{amsmath,amssymb,mathtools}
\usepackage[scr=rsfs,cal=euler]{mathalfa}
\usepackage{xfrac}
\usepackage[unicode,hidelinks,bookmarks=false]{hyperref}
\newcommand{\ie}{i.e.\ }
\newcommand{\cf}{cf.\ }
\newcommand{\resp}{respectively\ }
\newcommand{\Subsection}{\subsection}
\providecommand{\nobreakdash}{\nobreak\hbox{-}}
\setlength{\emergencystretch}{2em}
\multlinegap0pt
\usepackage{amssymb}
\usepackage{tikz}
\usepackage{xcolor}
\newtheorem{theorem}{Theorem}
\newcommand{\NN}[0]{\mathbb{N}}
\newcommand{\ludovic}[1]{\textcolor{blue}{#1}}
\newcommand{\quentin}[1]{\textcolor{green}{#1}}
\title{Ramsey-like theorems for separable permutations}
\author{Quentin Le Hou\'erou, Ludovic Patey}
\date{Source: arXiv:2507.07606v2 (CC BY 4.0)}
\begin{document}
\maketitle

\noindent\emph{Source and licence.} Ramsey-like theorems for separable permutations. Version arXiv:2507.07606v2 (CC BY 4.0), distributed by arXiv under the Creative Commons Attribution 4.0 International licence, http://creativecommons.org/licenses/by/4.0/, as recorded in the arXiv licence metadata for this exact version and shown on the abstract page https://arxiv.org/abs/2507.07606v2. Full licence notice: \texttt{licenses/arxiv-2507.07606-CC-BY-4.0.txt}.\par\smallskip

\noindent\textit{Editorial scope.} Counted unit: the article's theorem theorem (source0.tex lines 1543--1545). The article's complete original proof of it is delivered with it (source0.tex lines 1548--1589, one \texttt{proof} environment of 42 lines). No mathematics is written, repaired or re-derived: the statement and the proof are the article's own lines, in the article's own notation, and no retained line is substituted or re-flowed.  No further source-local context is needed: every cross-reference of every retained line resolves inside the two delivered spans.  Nothing was omitted from any retained span, so the package carries no omission record.  No retained line contains a citation, so the wrapper carries no bibliography and no bibliography entry.  No retained line contains a figure environment, an image-inclusion command or drawing code, so no image is delivered and the package declares no asset dependency.  Editorial formatting only, in addition: an AMS \texttt{amsart} preamble that restates the article's own declarations and macros used by the retained lines, namely the environments \texttt{theorem} on separate counters, together with the standard packages those lines need.  The retained environments print in this document as Theorem 1 in order of appearance, whereas the article prints the same items with the numbering of its own full document.  The complete native source of the article as distributed in the arXiv e-print for this exact version is bundled unchanged as \texttt{sources/181355/source0.tex}.
\begin{theorem}\label[theorem]{thm:linear-probabilistic-noncomputable}
    There exists a computable linear ordering $(\NN, \sqsubseteq)$ with no computable infinite monotone sequence, such that the measure of oracle computing an infinite increasing sequence is~1.
\end{theorem}

\begin{proof}
It is sufficient to find a computable linear order $(A, \sqsubseteq_A)$ for $A$ an infinite set, as well as a functional that outputs an infinite $\sqsubseteq_A$-increasing sequence on a set of measure $> \frac{1}{4}$.
For that, we will consider four functionals $\Gamma_{inc}$, $\Gamma_{dec}$, $\Delta_{inc}$ and $\Delta_{dec}$ with the following properties:
\begin{itemize}
    \item For $A = \{x_0, x_1, \dots \}$ a set and $e$ a Turing index, $\Gamma^A_{inc}(e)$ and $\Gamma^A_{dec}(e)$ will compute a linear order $(B,\sqsubseteq_B)$ for some $B \subseteq A$. If $A$ is infinite, then so will be $B$.
    
    \item For every sets $A$ and $C$, $\Delta^{C\oplus A}_{inc}(e)$ (resp $\Delta^{C\oplus A}_{dec}(e)$) will compute an increasing sequence of $\Gamma^A_{inc}(e)$ (resp $\Gamma^A_{dec}(e)$). This sequence will also be increasing for the standard ordering $<$ of the integers.
 
    If $A$ is infinite, then the measure of oracles $C$ such that $\Delta^{C\oplus A}_{inc}(e)$ (resp $\Delta^{C\oplus A}_{dec}(e)$) is an infinite increasing sequence will be at least $\Pi_{\ell > e} (1 - \frac{1}{2^{\ell}})$. Note that $\Pi_{\ell > 0} (1 - \frac{1}{2^{\ell}}) > \frac{1}{4}$.
    
    \item No $W_{e'} \cap B$ for $e' > e$ will be an infinite monotone sequence of $\Gamma^A_{inc}(e)$ or $\Gamma^A_{dec}(e)$. Furthermore, $W_{e} \cap B$ will not be an infinite monotone
    %\ludovic{increasing?}\quentin{Oui, c'est peut-être plus clair de dire increasing, mais la construction donne aussi le résultat pour les decreasing, vu qu'on $\Gammma_{dec}(e)$ fait appel à $\Gamma_{inc}(e)$}\ludovic{Effectivement, je ne l'ai compris qu'en lisant la suite de la construction.}\quentin{C'est peut être plus clair de ne pas parler des $e'$ ici et de juste mentionner increasing et decreasing}\ludovic{Oui, je pense}\ludovic{Finalement, non, je ne vois pas comment on peut juste mentionner le niveau~$e$ sans parler des $e' > e$. Sinon, on ne pourrait pas utiliser cette propriété en tant que blackbox pour obtenir la propriété globale, parce qu'en tant que blackbox, on ne sait pas que $\Gamma^A_{inc}(e)$ est obtenu à partir de $\Gamma^A_{dec}(e)$}
    sequence of $\Gamma^A_{dec}(e)$ and will not be an infinite decreasing sequence of $\Gamma^A_{inc}(e)$.
\end{itemize}

These four functionals will construct their ordering step by step, with the position of $x_s$ in the ordering (or the absence of $x_s$ in the ordering) being determined at stage $s$. $\Gamma_{inc}$ and $\Gamma_{dec}$ will call each other recursively as follows: For every infinite oracle~$A \subseteq \NN$, $\Gamma^A_{dec}(0)$ will call $\Gamma^A_{inc}(0)$, which will call $\Gamma^B_{dec}(1)$ for multiple~$B \subseteq A$, calling $\Gamma^B_{inc}(1)$, and so on... \\

Let $A = \{x_0, x_1, \dots \}$ be a set and $e$ be a Turing index. \\

\textbf{Definition of $\Gamma^A_{dec}(e)$.}
$\Gamma^A_{dec}(e)$ computes the order $(B, \sqsubseteq_B)$ defined as follows: $\Gamma^A_{dec}(e)$ is the same ordering as $\Gamma^A_{inc}(e)$, but, if at some stage $s$ we have $x_i \in W_e[s] \cap B$ for some $i < s$, then, all the $x_j$ for $j \geq s$ will be bigger than all the $x_i$ for $i < s$. The relative position of the $\{x_0, \dots, x_{s-1}\}$ and of the $\{x_{s}, \dots \}$ remain the one given by $\Gamma^{A}_{inc}(e)$.
This construction ensures that $W_e \cap B$ is not an infinite decreasing sequence of $\Gamma^A_{dec}(e)$, indeed, with this construction, there are only finitely many elements smaller than $x_i$ in the ordering $\Gamma^A_{dec}(e)$. \\

\textbf{Definition of $\Delta^{C \oplus A}_{dec}(e)$.}
Notice that any sequence that is both increasing for $\Gamma_{inc}^A(e)$ and for $<$ will also be an increasing sequence of $\Gamma_{dec}^A(e)$. Hence $\Delta^{C \oplus A}_{dec}(e)$ can be chosen to be the same as $\Delta^{C \oplus A}_{inc}(e)$ \\

\textbf{Definition of $\Gamma^A_{inc}(e)$.}
$\Gamma^A_{inc}(e)$ computes the order $(B, \sqsubseteq_B)$ defined as follows: at the first $2^{e+1}$ stages of the construction, we let $\{x_0, \dots, x_{2^{e+1}-1}\}$ be in $B$ with $x_0 \sqsubseteq_B \dots \sqsubseteq_B x_{2^{e+1} - 1}$, then, we partition the remaining elements $A \setminus \{x_0, \dots, x_{2^{e+1} - 1}\}$ into $2^{e+1}$ blocks $A_i = \{x_j \in A : j \geq 2^{e+1} \wedge j \equiv i \mbox{ mod } 2^{e+1}\}$ and let $(B_i, \sqsubseteq_{B_i})$ be the order computed by $\Gamma^{A_i}_{dec}(e+1)$ for every $i < 2^{e+1}$. The set $B$ will be a subset of $\{x_0, \dots, x_{2^{e+1} - 1}\} \cup \bigcup_{i < 2^{e+1}} B_i$, and the order $\sqsubseteq_B$ will coincide with the $\sqsubseteq_{B_i}$ on $B_i \cap B$ and will be such that $x_0 \sqsubseteq_B B_0 \cap B \sqsubseteq_B x_1 \sqsubseteq_B B_1 \cap B \sqsubseteq_B x_2 \sqsubseteq_B \dots \sqsubseteq_B B_{2^{e+1} - 1} \cap B$. To find out which elements of $\bigcup_{i < 2^{e+1}} B_i$ are in $B$, we proceed as follows: at every stage $s \geq 2^{e+1}$ of the construction, a single block $B_i$ will be disabled, which ensures that $x_s$ will be added to $B$ only if $x_s \not\in B_i$. Initially at stage $s_0 = 2^{e+1}$, $B_0$ is disabled, and if at some stage $s_1 \geq s_0$ an element $y_1 \in W_e[s_1]$ was already put in $B \cap B_{i_1}$ at a previous stage for some $0 < i_1 < 2^{e+1}$, then, we re-enable $B_0$ and disable $B_{i_1}$ until we found some stage $s_2$ and some element $y_2 \in W_e[s_2]$ that was already in $B \cap B_{i_2}$ for some $i_2 > i_1$, then we re-enable $B_{i_1}$ and disable $B_{i_2}$, and so on. % until we find some $s_3$....

This construction is well-defined. Indeed, to compute the position (or absence) of an element $x_s$ in the ordering, a finite amount of recursive calls are needed. This is due to the fact that $\Gamma^A_{inc}(e)$ immediately fixes the position of $x_0, \dots, x_{2^{e+1} - 1}$ without any recursive call, then, the position of the next $2^{e+2}$ elements only need the recursive call to $\Gamma^{A_i}_{dec}(e+1)$ which then call $\Gamma^{A_i}_{inc}(e+1)$, and so on for the next elements. %, which gives the position....

This definition of $\Gamma^A_{inc}(e)$ ensures that $W_e \cap B$ will not be an infinite increasing sequence. Indeed, otherwise, there would be some $i < 2^{e+1}$ such that after a certain stage $s$, all the elements of $W_e \cap B$ are in $B_i$ (and since the sequence is increasing, no elements of $W_e \cap B$ are in $B_j$ for $j > i$), this would disable $B_i$ with no reactivation for the rest of the construction, ensuring that $B_i \cap B$ is finite and contradicting our assumption. \\

\textbf{Definition of $\Delta^{C \oplus A}_{inc}(e)$.}
Notice that in the construction of the order $\Gamma^A_{inc}(e)$, exactly one of the blocks $B_i$ will end up disabled, thus, for all the other blocks $B_j$, an infinite increasing sequence of $(B_j, \sqsubseteq_{B_j})$ will output an infinite increasing sequence of $(B, \sqsubseteq_B)$ by removing the finite amount of elements that are not in $B_j \cap B$. Hence, $\Delta^{C \oplus A}_{inc}(e)$ can be defined to act as follows: using the first $e+1$ bits of $C$ it will pick one $j < 2^{e+1}$, then $x_j$ will be the smallest element of the increasing sequence outputted, and it will then recursively call $\Delta^{(C - (e + 1)) \oplus A_j}_{dec}(e+1)$ and remove the elements of the sequence that are not in $B_j \cap B$ to find the rest of the sequence.

$\Delta^{C \oplus A}_{inc}(e)$ will output an infinite increasing sequence if it never picks the wrong block that end up disabled during the recursive calls. This happens with probability $\Pi_{\ell > e} (1 - \frac{1}{2^{\ell}})$. \\

So far, we only proved that $\Gamma^{A}_{dec}(e)$ (resp $\Gamma^{A}_{inc}(e)$) ensures that $W_e \cap B$ will not be an infinite decreasing (resp increasing) sequence of the computed order. Due to the recursive nature of the functionals, we can go further than that and obtain the same property for every $e' > e$. Indeed, for every $e' > e$, all the elements ordered by $\Gamma^{A}_{inc}(e)$, except for a finite amount, have been ordered by one of the (finitely many) recursive call to some $\Gamma^{A_i}_{inc}(e')$ for some $A_i \subseteq A$. If $W_{e'} \cap B$ was an infinite increasing sequence, then this would yield an infinite decreasing subsequence for one of the $\Gamma^{A_i}_{inc}(e')$, contradicting what was proved. \\

Thus, taking $\Gamma_{dec}^{\NN}(0)$ to be the ordering and $\Delta^{(\cdot) \oplus \NN}_{dec}(0)$ to be the functional yields the desired result.
\end{proof}

\end{document}
